% Processed from chunks/03_cohen_newton.md.
% Bilingual LaTeX fragment. Compile with a Chinese-capable engine such as XeLaTeX/LuaLaTeX.
% Images are preserved with project-root-relative paths and require graphicx in the parent document.

\section*{Chapter 7: The Grand Design -- A New Physics\\第七章：宏伟构想：一种新物理学}
\begin{flushright}
I. Bernard Cohen
\end{flushright}

\noindent\textbf{English.} The publication of Isaac Newton's \emph{Principia} in 1687 was one of the most notable events in the whole history of physical science. In it one may find the culmination of thousands of years of striving to comprehend the system of the world, the principles of force and of motion, and the physics of bodies moving in different media. It is no small testimony to the vitality of Newton's scientific genius that although the physics of the \emph{Principia} has been altered, improved, and challenged ever since, we still set about solving most problems of celestial mechanics and the physics of gross bodies by proceeding essentially as Newton did some 300 years ago. Newtonian principles of celestial mechanics guide our artificial satellites, our space shuttles, and every spacecraft we launch to explore the vast reaches of our solar system. And if this is not enough to satisfy the canons of greatness, Newton was equally great as a pure mathematician. He invented the differential and integral calculus (produced simultaneously and independently by the German philosopher Gottfried Wilhelm Leibniz), which is the language of physics; he developed the binomial theorem and various properties of infinite series; and he laid the foundations for the calculus of variations.
\par\noindent\textbf{中文。} 艾萨克·牛顿的《自然哲学的数学原理》（\emph{Principia}）于1687年出版，这是整个物理科学史上最重要的事件之一。在这部著作中，人们可以看到数千年来理解世界体系、力与运动的原理，以及不同介质中物体运动之物理学的努力达到了顶点。此后，《原理》中的物理学不断被修改、完善并受到挑战；然而直到今天，我们处理天体力学和宏观物体物理学的大多数问题时，实质上仍沿着牛顿约三百年前开辟的道路前进，这充分显示了牛顿科学天才的生命力。牛顿的天体力学原理指引着人造卫星、航天飞机，以及我们发射去探索太阳系广阔空间的每一艘航天器。若这些还不足以满足伟大的标准，那么牛顿作为纯数学家同样伟大。他发明了微分和积分学（德国哲学家戈特弗里德·威廉·莱布尼茨也同时且独立地创立了它），而微积分正是物理学的语言；他发展了二项式定理和无穷级数的多种性质；并为变分法奠定了基础。
\medskip

\noindent\textbf{English.} In optics, Newton began the experimental study of the analysis and composition of light, showing that white light is a mixture of light of many colors, each having a characteristic index of refraction. Upon these researches have risen the science of spectroscopy and the methods of color analysis. Newton invented a reflecting telescope and so showed astronomers how to transcend the limitations of telescopes built of lenses. All in all, his was a fantastic scientific achievement--of a kind that has never been equaled and may never be equaled again.
\par\noindent\textbf{中文。} 在光学方面，牛顿开创了对光的分析与合成的实验研究，证明白光是许多种颜色光的混合，而每一种颜色都有其特定的折射率。在这些研究的基础上，光谱学和颜色分析方法发展起来。牛顿发明了反射望远镜，由此向天文学家展示了如何超越透镜望远镜的局限。总而言之，他的科学成就是惊人的，其性质前无古人，或许也后无来者。
\medskip

\noindent\textbf{English.} In this book we shall deal exclusively with Newton's system of dynamics and gravitation, the central problems for which the preceding chapters have been a preparation. If you have read them carefully, you have in mind all but one of the major ingredients requisite to an understanding of the Newtonian system of the world. But even if that one were to be given--the analysis of uniform circular motion--the guiding hand of Newton would still be required to put the ingredients together. It took genius to supply the new concept of universal gravitation. Let us see what Newton actually did.
\par\noindent\textbf{中文。} 本书将专门讨论牛顿的动力学与引力体系，前面各章正是为这些核心问题所作的准备。如果你已仔细读过前文，那么理解牛顿世界体系所需的主要要素，除了一个之外，都已在你心中具备。可是即便再给出那一个要素，即匀速圆周运动的分析，仍然需要牛顿的引导之手将这些要素结合起来。提出万有引力这一新概念，需要天才。让我们看看牛顿实际做了什么。
\medskip

\noindent\textbf{English.} First of all, it must be understood that Galileo himself never attempted to display any scheme of forces that would account for the movement of the planets, or of their satellites. As for Copernicus, the \emph{De revolutionibus} contains no important insight into a celestial mechanics. Kepler had tried to supply a celestial mechanism, but the result was never a very happy one. He held that the \emph{anima motrix} emanating from the sun would cause planets to revolve about the sun in circles. He further supposed that magnetic interactions of the sun and a planet would shift the planet during an otherwise circular revolution into an elliptical orbit. Others who contemplated the problems of planetary motion proposed systems of mechanics containing certain features that were later to appear in Newtonian dynamics. One of these was Robert Hooke, who quite understandably thought that Newton should have given him more credit than a mere passing reference for having anticipated parts of the laws of dynamics and gravitation.
\par\noindent\textbf{中文。} 首先必须明白，伽利略本人从未试图提出一套力的方案，以解释行星或其卫星的运动。至于哥白尼，《天体运行论》中并没有关于天体力学的重要洞见。开普勒曾试图提供一种天体机制，但结果并不十分成功。他认为从太阳发出的\emph{anima motrix}（运动灵魂）会使行星围绕太阳作圆周运动。他还进一步设想，太阳与行星之间的磁相互作用，会把本来作圆周运行的行星推移到椭圆轨道上。其他思考行星运动问题的人也提出过一些力学体系，其中含有后来会出现在牛顿动力学中的若干特征。罗伯特·胡克就是其中之一；他认为自己预见了动力学与引力定律的一部分，牛顿本应给予他比一句顺带提及更多的承认，这种想法是可以理解的。
\medskip

\subsection*{Newtonian Anticipations\\牛顿思想的先声}

\noindent\textbf{English.} The climactic chapter in the discovery of the mechanics of the universe starts with a pretty story. By the third quarter of the seventeenth century, a group of men had become so eager to advance the new mathematical experimental sciences that they banded together to perform experiments in concert, to present problems for solution to one another, and to report on their own researches and on those of others as revealed by correspondence, books, and pamphlets. Thus it came about that Robert Hooke, Edmond Halley, and Sir Christopher Wren, England's foremost architect, met to discuss the question, Under what law of force would a planet follow an elliptical orbit? From Kepler's laws--especially the third or harmonic law, but also the second or law of areas--it was clear that the sun somehow or other must control or at least affect the motion of a planet in accordance with the relative proximity of the planet to the sun. Even if the particular mechanisms proposed by Kepler (an \emph{anima motrix} and a magnetic force) had to be rejected, there could be no doubt that some kind of planet-sun interaction keeps the planets in their courses. Furthermore, a more acute intuition than Kepler's would sense that any force emanating from the sun must spread out in all directions from that body, presumably diminishing according to the inverse of the square of its distance from the sun--as the intensity of light diminishes in relation to distance. But to say this much is a very different thing from proving it mathematically. For to prove it would require a complete physics with mathematical methods for solving all the attendant and consequent problems. When Newton declined to credit authors who tossed off general statements without being able to prove them mathematically or fit them into a valid framework of dynamics, he was quite justified in saying, as he did of Hooke's claims: ``Now is not this very fine? Mathematicians that find out, settle, and do all the business must content themselves with being nothing but dry calculators and drudges; and another, that does nothing but pretend and grasp at all things, must carry away all the invention, as well of those that were to follow him as of those that went before.'' (See, further, Supplement 11.)
\par\noindent\textbf{中文。} 宇宙力学发现史中最高潮的一章，是从一个颇动人的故事开始的。到17世纪第三个四分之一时期，一批人已急切地想推进新的数学实验科学，他们结成团体，共同做实验，彼此提出待解的问题，并通过通信、书籍和小册子报告自己的研究以及他人的研究。于是，罗伯特·胡克、埃德蒙·哈雷和英国首屈一指的建筑师克里斯托弗·雷恩爵士聚在一起，讨论这样一个问题：在什么样的力的定律下，行星会沿椭圆轨道运行？根据开普勒定律，尤其是第三定律即调和定律，同时也包括第二定律即面积定律，显然太阳必定以某种方式按照行星与太阳的相对距离来控制，或至少影响行星的运动。即使必须抛弃开普勒所提出的具体机制（运动灵魂和磁力），也毫无疑问，某种行星--太阳相互作用使行星保持在各自的轨道上。此外，比开普勒更敏锐的直觉会意识到，从太阳发出的任何力都必定从太阳向四面八方扩散，并且大概按距太阳距离平方的倒数减弱，正如光强随距离而减弱一样。但是，说到这一点与用数学证明这一点，是完全不同的两回事。要证明它，就需要一套完整的物理学，以及解决所有相关和后续问题的数学方法。因此，当有些作者只能随口提出一般性断言，却既不能作数学证明，也不能把这些断言纳入有效的动力学框架时，牛顿拒绝给予他们荣誉是有充分理由的。对于胡克的主张，牛顿曾说：“这岂不是妙极了吗？数学家们发现、确定并做完全部工作，却只能满足于不过是枯燥的计算者和苦工；而另一个人什么也不做，只是自称什么都懂、什么都要抓到手，却必须把所有发明都归于自己，既包括后来者的，也包括前人已有的。”（另见补编11。）
\medskip

\noindent\textbf{English.} In any event, by January 1684 Halley had concluded that the force acting on planets to keep them in their orbits ``decreased in the proportion of the squares of the distances reciprocally,''
\[
F \propto \frac{1}{D^{2}},
\]
but he was not able to deduce from that hypothesis the observed motions of the celestial bodies. When Wren and Hooke met later in the month, they agreed with Halley's supposition of a solar force. Hooke boasted ``that upon that principle all the laws of the celestial motions were to be [i.e., could be] demonstrated, and that he himself had done it.'' But despite repeated urgings and Wren's offer of a considerable monetary prize, Hooke did not--and presumably could not--produce a solution. Six months later, in August 1684, Halley decided to go to Cambridge to consult Isaac Newton. On his arrival he learned the ``good news'' that Newton ``had brought this demonstration to perfection.'' Here is De Moivre's almost contemporaneous account of that visit:
\par\noindent\textbf{中文。} 无论如何，到1684年1月，哈雷已经得出结论：作用在行星上、使其保持在轨道中的力，“按距离平方的倒数比例减小”，
\[
F \propto \frac{1}{D^{2}},
\]
但他无法从这一假设推出天体的观测运动。那个月稍晚些时候，雷恩和胡克会面时，他们同意哈雷关于太阳力的设想。胡克夸口说，“依照那个原理，天体运动的一切定律都将[也就是能够]得到证明，而且他本人已经做到了”。可是，尽管众人一再催促，雷恩还提出一笔可观的奖金，胡克并没有拿出解答，而且大概也拿不出来。六个月后，即1684年8月，哈雷决定去剑桥请教艾萨克·牛顿。他到达后得知了一个“好消息”：牛顿“已经把这个证明完成得尽善尽美”。以下是德·莫阿弗对那次访问所作的几乎同时代的叙述：
\medskip

\begin{quote}
\noindent\textbf{English.} After they had been some time together, the Dr. [Halley] asked him what he thought the curve would be that would be described by the planets supposing the force of attraction towards the sun to be reciprocal to the square of their distance from it. Sir Isaac replied immediately that it would be an ellipsis. The Doctor, struck with joy and amazement, asked him how he knew it. Why, saith he, I have calculated it. Whereupon Dr. Halley asked him for his calculation without any further delay. Sir Isaac looked among his papers but could not find it, but he promised him to renew it and then to send it him. Sir Isaac, in order to make good his promise, fell to work again, but he could not come to that conclusion which he thought he had before examined with care. However, he attempted a new way which, though longer than the first, brought him again to his former conclusion. Then he examined carefully what might be the reason why the calculation he had undertaken before did not prove right, and he found that, having drawn an ellipsis coarsely with his own hand, he had drawn the two axes of the curve, instead of drawing two diameters somewhat inclined to one another, whereby he might have fixed his imagination to any two conjugate diameters, which was requisite he should do. That being perceived, he made both his calculations agree together.
\par\noindent\textbf{中文。} 两人相处了一段时间之后，博士[哈雷]问他：假定指向太阳的吸引力与行星距太阳距离的平方成反比，那么行星将描画出什么曲线？艾萨克爵士立刻回答说，会是一个椭圆。博士又惊又喜，问他是怎么知道的。他说，哦，我算过。于是哈雷博士不再迟疑，请他拿出计算。艾萨克爵士在自己的论文中寻找，却找不到；但他答应重新作一次，再寄给哈雷。为了履行诺言，艾萨克爵士重新动手，可是他无法得到自己以为先前已经仔细考察过的结论。不过，他尝试了另一种方法；这种方法虽比第一种更长，却又把他带回到原先的结论。随后他仔细检查，想找出先前所作计算为何不正确；他发现，自己用手粗略画椭圆时，画的是曲线的两条轴，而不是两条彼此略有倾斜的直径。要使想象固定在任意两条共轭直径上，后者才是必须画的。意识到这一点后，他使两种计算彼此一致了。
\end{quote}

\noindent\textbf{English.} Spurred on by Halley's visit, Newton resumed work on a subject that had commanded his attention in his twenties when he had laid the foundations of his other great scientific discoveries: the nature of white light and color and the differential and integral calculus. He now put his investigations in order, made great progress, and in the fall term of the year, discussed his research in a series of lectures on dynamics that he gave at Cambridge University, as required by his professorship. Eventually, with Halley's encouragement, a draft of these lectures, \emph{De motu corporum}, grew into one of the greatest and most influential books any man has yet conceived. Many a scientist has echoed the sentiment that Halley expressed in the ode he wrote as a preface to Newton's \emph{Principia} (or, to give Newton's masterpiece its full title, \emph{Philosophiae naturalis principia mathematica}, \emph{Mathematical Principles of Natural Philosophy}, London, 1687):
\par\noindent\textbf{中文。} 在哈雷来访的推动下，牛顿重新研究起一个在他二十多岁时就曾吸引其注意力的课题；那时他已经为另外两项伟大的科学发现奠定了基础，即白光与颜色的本性，以及微分和积分学。现在他整理自己的研究，取得重大进展，并在当年秋季学期按教授职务的要求，在剑桥大学作了一系列关于动力学的讲座，讨论这些研究。最终，在哈雷的鼓励下，这些讲稿的草稿《物体运动论》（\emph{De motu corporum}）发展成为人类迄今构思出的最伟大、最有影响力的著作之一。许多科学家都呼应过哈雷在为牛顿《原理》所写序诗中表达的感情（若写出牛顿这部杰作的全名，则是\emph{Philosophiae naturalis principia mathematica}，《自然哲学的数学原理》，伦敦，1687年）：
\medskip

\begin{verse}
Then ye who now on heavenly nectar fare,\\
Come celebrate with me in song the name\\
Of Newton, to the Muses dear; for he\\
Unlocked the hidden treasuries of Truth:\\
So richly through his mind had Phoebus cast\\
The radiance of his own divinity.\\
Nearer the gods no mortal may approach.
\end{verse}

\begin{verse}
那么，你们这些如今享用天上琼浆的人，\\
请与我同歌牛顿之名，\\
他为缪斯所珍爱；因为他\\
开启了真理深藏的宝库：\\
福玻斯曾如此丰沛地\\
把自身神性的光辉投射进他的心灵。\\
凡人再也不能更接近诸神。
\end{verse}

\subsection*{The Principia\\《原理》}

\noindent\textbf{English.} The \emph{Principia} is divided into three parts or books; we shall concentrate on the first and third. In Book One Newton develops the general principles of the dynamics of moving bodies, and in Book Three he applies the principles to the mechanism of the universe. Book Two deals with a facet of fluid mechanics, the theory of waves, and other aspects of physics.
\par\noindent\textbf{中文。} 《原理》分为三部分，或三卷；我们将集中讨论第一卷和第三卷。在第一卷中，牛顿发展了运动物体的动力学一般原理；在第三卷中，他把这些原理应用于宇宙的机制。第二卷处理流体力学的一个方面、波的理论以及物理学的其他方面。
\medskip

\noindent\textbf{English.} In Book One, following the preface, a set of definitions, and a discussion of the nature of time and space, Newton presented the ``axioms, or laws of motion.''
\par\noindent\textbf{中文。} 在第一卷中，牛顿在序言、一组定义以及关于时间与空间性质的讨论之后，提出了“公理，或运动定律”。
\medskip

\subsubsection*{Law I\\定律一}
\noindent\textbf{English.} Every body perseveres in its state of being at rest or of moving uniformly straight forward, except insofar as it is compelled to change its state by forces impressed upon it.
\par\noindent\textbf{中文。} 每一个物体都会保持其静止状态，或沿直线作匀速运动的状态，除非它被施加在其上的力迫使改变这种状态。
\medskip

\subsubsection*{Law II\\定律二}
\noindent\textbf{English.} A change in motion is proportional to the motive force impressed and takes place in the direction of the straight line along which that force is impressed. [See Supplementary Note on p. 184.]
\par\noindent\textbf{中文。} 运动的改变与所施加的动力成正比，并沿着施加该力的直线方向发生。[见第184页补注。]
\medskip

\begin{center}
\includegraphics[width=\linewidth]{chunks/images/0f043290c0aee6258d3b9dc29864191fc41e15048c067bc65b960e2596f4c37d.jpg}
\end{center}

\noindent\textbf{Plate VI.} A ball fired with a spring gun in the smokestack of a moving toy locomotive describes a parabola and lands on the locomotive instead of going straight up and down as it does when the locomotive is standing still. These stroboscopic pictures, with exposures at intervals of one-thirtieth of a second, vividly illustrate one of Galileo's arguments on the behavior of falling bodies and settle the ancient debate about bodies dropped from the masts of moving ships. If the speed of the locomotive were absolutely uniform and if the ball met no air resistance, it would land in the smokestack. (In fact, even under the imperfect conditions of the experiment, the ball hits the smokestack more often than not.) Note that the ball attains the same height whether the locomotive is at rest or moving. Notice, too, that in the picture where the locomotive is standing still, the distances traveled by the ball in the intervals of exposure correspond almost exactly. On the ascent, gravity slows it down; on the descent, gravity speeds it up. Photographs by Berenice Abbott.
\par\noindent\textbf{图版VI。} 在一辆运动中的玩具机车烟囱里，用弹簧枪射出一个小球；小球描出一条抛物线，并落回机车上，而不是像机车静止时那样笔直上升又下降。这些频闪照片以三十分之一秒为曝光间隔，生动说明了伽利略关于落体行为的一个论证，也解决了关于从运动船只桅杆上落下物体的古老争论。如果机车速度绝对均匀，并且小球不受空气阻力，小球就会落回烟囱中。（事实上，即使在实验条件并不完美的情况下，小球击中烟囱的次数也多于未击中的次数。）注意，无论机车静止还是运动，小球达到的高度都相同。还应注意，在机车静止的照片中，小球在各曝光间隔内通过的距离几乎正好对应。上升时，重力使它减速；下降时，重力使它加速。照片由贝伦妮丝·阿博特拍摄。
\medskip

\noindent\textbf{English.} Observe that if a body is in uniform motion in a straight line, a force at right angles to the direction of motion of the body will not affect the forward motion. This follows from the fact that the acceleration is always in the same direction as the force producing it, so that the acceleration in this case is at right angles to the direction of motion. Thus in the toy train experiment of Chapter 5, the chief force acting is the downward force of gravity, producing a vertical acceleration. The ball, whether moving forward or at rest, is thus caused to slow down in its upward motion until it comes to rest, and then be speeded up or accelerated on the way down.
\par\noindent\textbf{中文。} 请注意，如果一个物体正沿直线作匀速运动，那么与该物体运动方向成直角的力不会影响其前进运动。这是因为加速度总是与产生它的力同方向，因此在这种情况下，加速度与运动方向成直角。所以，在第五章的玩具火车实验中，起主要作用的力是向下的重力，它产生竖直加速度。于是，无论小球本来是否向前运动，它在向上运动时都会减速，直到停止；随后在下降途中又会加速。
\medskip

\noindent\textbf{English.} A comparison of the two sets of photographs (p. 83) shows that the upward and downward motions are exactly the same whether the train is at rest or in uniform motion. In the forward direction there is no effect of weight or gravity, since this acts only in a downward direction. The only force in the forward or horizontal direction is the small amount of air friction, which is almost negligible; so one may say that in the horizontal direction there is no force acting. According to Newton's first law of motion, the ball will continue to move in the forward direction with uniform motion in a straight line just as the train does--a fact you can check by inspecting the photograph. The ball remains above the locomotive whether the train is at rest or in uniform motion in a straight line. This law of motion is sometimes called the principle of inertia, and the property that material bodies have of continuing in a state of rest or of uniform motion in a straight line is sometimes known as the bodies' inertia.\textsuperscript{1}
\par\noindent\textbf{中文。} 比较两组照片（第83页）可见，无论火车静止还是作匀速运动，小球的上升和下降运动都完全相同。在前进方向上，重量或重力不起作用，因为重力只沿向下方向作用。前进方向或水平方向上唯一的力，是很小的一点空气摩擦，几乎可以忽略；因此可以说，水平方向上没有力在作用。根据牛顿第一运动定律，小球将像火车一样继续沿前进方向作直线匀速运动；这一点可以通过观察照片来核对。无论火车静止还是沿直线匀速运动，小球都保持在机车上方。这条运动定律有时称为惯性原理，而物质物体继续保持静止状态或直线匀速运动状态的性质，有时称为物体的惯性。\textsuperscript{1}
\medskip

\noindent\textbf{English.} Newton illustrated Law I by reference to projectiles that continue in their forward motions ``so far as they are not retarded by the resistance of the air, or impelled downward by the force of gravity,'' and he referred also to ``the greater bodies of planets and comets.'' (On the inertial aspect of the motion of ``greater bodies'' such as ``planets and comets,'' see Supplement 12.) At this one stroke Newton postulated the opposite view of Aristotelian physics. In the latter, no celestial body could move uniformly in a straight line in the absence of a force, because this would be a ``violent'' motion and so contrary to its nature. Nor could a terrestrial object, as we have seen, move along its ``natural'' straight line without an external mover or an internal motive force. Newton, presenting a physics that applies simultaneously to both terrestrial and celestial objects, stated that in the absence of a force bodies do not necessarily stand still or come to rest as Aristotle supposed, but they may move at constant rectilinear speed. This ``indifference'' of all sorts of bodies to rest or uniform straight-line motion in the absence of a force clearly is an advanced form of Galileo's statement in his book on sunspots (p. 88), the difference being that in that work Galileo was writing about uniform motion along a great spherical surface concentric with the earth.
\par\noindent\textbf{中文。} 牛顿用抛射体来说明定律一：只要抛射体“未被空气阻力所阻滞，或未被重力向下推动”，它们就会继续向前运动；他还提到“行星和彗星这些较大的物体”。（关于“行星和彗星”等“较大物体”运动中的惯性方面，见补编12。）仅此一举，牛顿就提出了与亚里士多德物理学相反的观点。在亚里士多德物理学中，没有力时任何天体都不可能沿直线匀速运动，因为这会是“强迫”运动，因而违背其本性。同样，正如我们已经看到的，地上物体如果没有外在推动者或内在动力，也不能沿其“自然”直线运动。牛顿提出了一种同时适用于地上物体和天上物体的物理学，指出在没有力的情况下，物体并不一定像亚里士多德所设想的那样静止不动或趋于静止，而是可以以恒定速度作直线运动。在没有力时，各类物体对于静止和直线匀速运动表现出的这种“无差别”，显然是伽利略在其太阳黑子著作（第88页）中说法的进一步形式；不同的是，伽利略在那里谈的是沿着与地球同心的巨大球面作匀速运动。
\medskip

\noindent\textbf{English.} Newton said of the laws of motion that they were ``such principles as have been received by mathematicians, and \ldots{} confirmed by [an] abundance of experiments. By the first two Laws and the first two Corollaries, Galileo discovered that the descent of bodies varies as the square of the time and that the motion of projectiles is in the curve of a parabola, experience agreeing with both, unless so far as these motions are a little retarded by the resistance of the air.'' The ``two Corollaries'' deal with methods used by Galileo and many of his predecessors to combine two different forces or two independent motions. Fifty years after the publication of Galileo's \emph{Two New Sciences} it was difficult for Newton, who had already established an inertial physics, to conceive that Galileo could have come as close as he had to the concept of inertia without having taken full leave of circularity and having known the true principle of linear inertia.
\par\noindent\textbf{中文。} 牛顿论及运动定律时说，它们是“数学家们已经接受，并且……由大量实验所证实的原理。凭借前两条定律和前两个推论，伽利略发现物体下落与时间的平方成比例，抛射体的运动沿抛物线进行；经验与这两点一致，只是这些运动会因空气阻力而略有迟滞”。所谓“两个推论”，涉及伽利略及其许多前人用来合成两个不同的力或两个独立运动的方法。在伽利略《两门新科学》出版五十年后，对于已经建立惯性物理学的牛顿来说，很难想象伽利略竟能如此接近惯性概念，却没有完全摆脱圆形运动，也没有认识真正的线性惯性原理。
\medskip

\noindent\textbf{English.} Newton was being very generous to Galileo because, however it may be argued that Galileo ``really did'' have the law of inertia or Newton's Law I, a great stretch of the imagination is required to assign any credit to Galileo for Law II. This law has two parts. In the second half of Newton's statement of Law II, the ``change in motion'' produced by an ``impressed'' or ``motive'' force--whether that is a change in the speed with which a body moves or a change in the direction in which it is moving--is said to be ``in the direction of the straight line along which that force is impressed.'' This much is certainly implied in Galileo's analysis of projectile motion because Galileo assumed that in the forward direction there is no acceleration because there is no horizontal force, except the negligible action of air friction; but in the vertical direction there is an acceleration or continual increase of downward speed, because of the downward-acting weight force. But the first part of Law II--that the change in the magnitude of the motion is related to the motive force--is something else again; only a Newton could have seen it in Galileo's studies of falling bodies. This part of the law says that if an object were to be acted on first by one force $F_{1}$ and then by some other force $F_{2}$, the accelerations or changes in speed produced, $A_{1}$ and $A_{2}$, would be proportional to the forces, or that
\[
\frac{F_{1}}{F_{2}} = \frac{A_{1}}{A_{2}}, \qquad \mathrm{or} \qquad
\frac{F_{1}}{A_{1}} = \frac{F_{2}}{A_{2}}.
\]
\par\noindent\textbf{中文。} 牛顿对伽利略已经十分慷慨，因为无论人们如何主张伽利略“确实”已经具有惯性定律或牛顿定律一的思想，要把定律二的功劳也归给伽利略，都需要极大的想象力。这条定律有两个部分。在牛顿对定律二的表述后半部分中，由“施加的”力或“动力”所产生的“运动的改变”——无论是物体运动速度的改变，还是运动方向的改变——都被说成“沿着施加该力的直线方向”发生。这一点当然隐含在伽利略对抛射体运动的分析中，因为伽利略假定，在前进方向上除几乎可忽略的空气摩擦外没有水平力，所以没有加速度；而在竖直方向上，由于有向下作用的重量力，存在加速度，或者说向下速度不断增加。但是定律二的前半部分，即运动量大小的改变与动力有关，却是另一回事；只有牛顿才能在伽利略的落体研究中看出这一点。定律的这一部分说，如果一个物体先受一个力$F_{1}$作用，然后又受另一个力$F_{2}$作用，那么所产生的加速度或速度改变$A_{1}$和$A_{2}$将与这些力成正比，或者说
\[
\frac{F_{1}}{F_{2}} = \frac{A_{1}}{A_{2}}, \qquad \hbox{或} \qquad
\frac{F_{1}}{A_{1}} = \frac{F_{2}}{A_{2}}.
\]
\medskip

\noindent\textbf{English.} But in analyzing falling, Galileo was dealing with a situation in which only one force acted on each body, its weight $W$, and the acceleration it produced was $g$, the acceleration of a freely falling body. (For the two forms of Newton's Law II, see p. 184.)
\par\noindent\textbf{中文。} 可是在分析下落时，伽利略处理的是每个物体只受一个力作用的情形，即它的重量$W$；这个力所产生的加速度是$g$，即自由落体的加速度。（关于牛顿定律二的两种形式，见第184页。）
\medskip

\noindent\textbf{English.} Where Aristotle had said that a given force gives an object a certain characteristic speed, Newton now said that a given force always produces in that body a definite acceleration $A$. To find the speed $V$, we must know how long a time $T$ the force has acted, or how long the object has been accelerated, so that Galileo's law
\[
V = AT
\]
may be applied.
\par\noindent\textbf{中文。} 亚里士多德曾说，给定的力会给予物体某种特定速度；而牛顿现在说，给定的力总是在该物体中产生确定的加速度$A$。要找到速度$V$，我们必须知道力作用了多长时间$T$，也就是物体被加速了多长时间，这样才可以应用伽利略定律
\[
V = AT.
\]
\medskip

\noindent\textbf{English.} At this point let us try a thought experiment, in which we assume we have two cubes of aluminum, one just twice the volume of the other. (Incidentally, to ``duplicate'' a cube--or make a cube having exactly twice the volume as some given cube--is as impossible within the framework of Euclidean geometry as to trisect an angle or to square a circle.) We now subject the smaller cube to a series of forces $F_{1}, F_{2}, F_{3}, \ldots$ and determine the corresponding accelerations $A_{1}, A_{2}, A_{3}, \ldots$. In accordance with Law II, we would find that there is a certain constant value of the ratio of force to acceleration
\[
\frac{F_{1}}{A_{1}} = \frac{F_{2}}{A_{2}} = \frac{F_{3}}{A_{3}} = \dots = m_{s},
\]
which for this object we may call $m_s$.
\par\noindent\textbf{中文。} 此时让我们做一个思想实验。假定我们有两个铝立方体，其中一个的体积正好是另一个的两倍。（顺便说一句，在欧几里得几何框架内，“倍立方”——也就是作出一个体积正好为某给定立方体两倍的立方体——与三等分任意角或化圆为方一样不可能。）现在我们让较小的立方体依次受到一系列力$F_{1}, F_{2}, F_{3}, \ldots$，并测定相应的加速度$A_{1}, A_{2}, A_{3}, \ldots$。根据定律二，我们会发现力与加速度之比有某个恒定值：
\[
\frac{F_{1}}{A_{1}} = \frac{F_{2}}{A_{2}} = \frac{F_{3}}{A_{3}} = \dots = m_{s},
\]
对于这个物体，我们可以把它称为$m_s$。
\medskip

\noindent\textbf{English.} We now repeat the operations with the larger cube and find that the same set of forces $F_1, F_2, F_3, \ldots$ respectively produces another set of accelerations $a_1, a_2, a_3, \ldots$. In accordance with Newton's second law, the force-acceleration ratio is again a constant which for this object we may call $m_l$:
\[
\frac{F_{1}}{a_{1}} = \frac{F_{2}}{a_{2}} = \frac{F_{3}}{a_{3}} = \dots = m_{l}.
\]
\par\noindent\textbf{中文。} 现在我们对较大的立方体重复这些操作，发现同一组力$F_1, F_2, F_3, \ldots$分别产生另一组加速度$a_1, a_2, a_3, \ldots$。根据牛顿第二定律，力与加速度之比再次是一个常数；对于这个物体，我们可以称之为$m_l$：
\[
\frac{F_{1}}{a_{1}} = \frac{F_{2}}{a_{2}} = \frac{F_{3}}{a_{3}} = \dots = m_{l}.
\]
\medskip

\noindent\textbf{English.} For the larger object the constant proves to be just twice as large as the constant obtained for the smaller one and, in general, so long as we deal with a single variety of matter like pure aluminum, this constant is proportional to the volume and so is a measure of the amount of aluminum in any sample. This particular constant is a measure of an object's resistance to acceleration, or a measure of the tendency of that object to stay as it is--either at rest, or in motion in a straight line. For observe that $m_l$ was twice $m_s$; to give both objects the same acceleration or change in motion the force required for the larger object is just twice what it must be for the smaller. The tendency of any object to continue in its state of motion (at constant speed in a straight line) or its state of rest is called its inertia; hence, Newton's Law I is also called the principle of inertia. The constant determined by finding the constant force-acceleration ratio for any given body may thus be called the body's inertia. But for our aluminum blocks this same constant is also a measure of the ``quantity of matter'' in the object, which is called its mass. We now make precise the condition that two objects of different material--say one of brass and the other of wood--shall have the same ``quantity of matter'': it is that they have the same mass as determined by the force-acceleration ratio, or the same inertia.
\par\noindent\textbf{中文。} 对较大的物体而言，这个常数恰好是较小物体所得常数的两倍；一般地说，只要我们处理的是像纯铝这样单一类型的物质，这个常数就与体积成正比，因此也是任何样品中铝的多少的量度。这个特定常数是物体抵抗加速的量度，或者说是物体保持原状的倾向的量度——无论原状是静止，还是沿直线运动。请注意，$m_l$是$m_s$的两倍；若要使两个物体获得相同的加速度或运动改变，较大物体所需的力正好是较小物体的两倍。任何物体继续保持其运动状态（以恒定速度沿直线运动）或静止状态的倾向称为它的惯性；因此，牛顿定律一也称为惯性原理。对于任一给定物体，通过求出恒定的力--加速度比所确定的常数，也可称为该物体的惯性。但对我们的铝块来说，这同一个常数也是物体中“物质多少”的量度，而这种量度称为它的质量。现在我们可以精确说明，不同材料的两个物体——比如一个黄铜物体和一个木头物体——具有相同“物质多少”的条件是什么：它们必须具有由力--加速度比确定的相同质量，或者说具有相同惯性。
\medskip

\noindent\textbf{English.} In ordinary life, we do not compare the ``quantity of matter'' in objects in terms of their inertias, but in terms of their weight. Newtonian physics makes it clear why we can, and through its clarification we are able to understand why at any place on the earth two unequal weights in a vacuum fall at the same rate. But we may observe that in at least one common situation we always compare the inertias of objects rather than their weights. This happens when a person hefts two objects to find which is heavier, or has the greater mass. He does not hold them out to see which pulls down more on his arm; instead, he moves them up and down to find which is easier to move. In this way he determines which has the greater resistance to a change in its state of motion in a straight line or of rest--that is, which has the greater inertia. (On Newton's concept of inertia, see Supplement 15.)
\par\noindent\textbf{中文。} 在日常生活中，我们并不根据物体的惯性来比较其中的“物质多少”，而是根据其重量来比较。牛顿物理学说明了我们为什么可以这样做；通过这种说明，我们也能够理解为什么在地球上任一地点，真空中两个重量不同的物体会以相同速率下落。不过我们可以注意到，至少在一种常见情形中，我们总是在比较物体的惯性而不是重量。这发生在一个人掂量两个物体，以判断哪一个更重或质量更大时。他并不是把它们伸出去，看哪一个把手臂向下拉得更多；相反，他上下移动它们，看哪一个更容易移动。通过这种方式，他判断哪一个对其直线运动状态或静止状态的改变有更大的抵抗，也就是哪一个有更大的惯性。（关于牛顿的惯性概念，见补编15。）
\medskip

\subsection*{Final Formulation of the Law of Inertia\\惯性定律的最终表述}

\noindent\textbf{English.} At one point in his \emph{Discourses and Demonstrations Concerning Two New Sciences}, Galileo imagined a ball to be rolling along a plane and noted that ``equable motion on this plane would be perpetual if the plane were of infinite extent.'' A plane without limit is all right for a pure mathematician, who is a Platonist in any case. But Galileo was a man who combined just such a Platonism with a concern for applications to the real world of sensory experience. In the \emph{Two New Sciences}, Galileo was not interested only in abstractions as such, but in the analysis of real motions on or near the earth. So we understand that having talked about a plane without limit, he does not continue with such a fancy, but asks what would happen on such a plane if it were a real earthly plane, which for him means that it is ``ended, and [situated] on high.'' The ball, in the real world of physics, falls off the plane and begins to fall to the ground. In this case,
\par\noindent\textbf{中文。} 在《两门新科学的对话与数学证明》的某处，伽利略想象一个球沿平面滚动，并指出：“如果这个平面无限延伸，其上的匀速运动就会永远持续。”一个没有边界的平面对于纯数学家来说当然可以接受，因为纯数学家无论如何都是柏拉图主义者。但伽利略把这种柏拉图主义与对感官经验中现实世界之应用的关切结合在一起。在《两门新科学》中，伽利略并不只关心抽象本身，而是关心地球上或地球附近真实运动的分析。因此我们理解，在谈到一个无边界平面之后，他并没有继续这种想象，而是问：如果这样的平面是一个真实的地上平面，会发生什么？对他来说，这意味着该平面是“有尽头的，并且[位于]高处”。在物理学的现实世界中，球会滚出平面，开始向地面下落。在这种情况下，
\medskip

\begin{quote}
\noindent\textbf{English.} the movable (which I conceive of as being endowed with heaviness), driven to the end of this plane and going on further, adds on to its previous equable and indelible motion that downward tendency which it has from its own heaviness. Thus there emerges a certain motion, compounded from equable horizontal and from naturally accelerated downward [motion], which I call ``projection.''
\par\noindent\textbf{中文。} 可动体（我设想它具有重性）被推到这个平面的尽头并继续前进时，会在其先前均匀且不可磨灭的运动之上，增添由自身重性产生的向下倾向。于是便出现一种运动，它由均匀的水平运动和自然加速的向下[运动]复合而成，我称之为“抛射”。
\end{quote}

\noindent\textbf{English.} Unlike Galileo, Newton made a clear separation between the world of abstract mathematics and the world of physics, which he still called philosophy. Thus the \emph{Principia} included both ``mathematical principles'' as such and those that could be applied in ``natural philosophy,'' but Galileo's \emph{Two New Sciences} included only those mathematical conditions exemplified in nature. For instance, Newton plainly knew that the attractive force exerted by the sun on a planet varies as the inverse square of the distance,
\[
F \propto \frac{1}{D^{2}},
\]
but in Book One of the \emph{Principia} he explored the consequences not only of this particular force but of others with quite different dependence on the distance, including
\[
F \propto D.
\]
\par\noindent\textbf{中文。} 与伽利略不同，牛顿在抽象数学世界和物理世界之间作出了清楚区分；他仍把物理世界称为哲学。因此，《原理》既包括作为数学本身的“数学原理”，也包括可以应用于“自然哲学”的原理；而伽利略的《两门新科学》只包括自然中所体现的那些数学条件。例如，牛顿清楚知道太阳对行星施加的吸引力随距离平方的倒数而变化，
\[
F \propto \frac{1}{D^{2}},
\]
但在《原理》第一卷中，他不仅考察了这种特定力的结果，也考察了其他与距离有截然不同依赖关系的力的结果，包括
\[
F \propto D.
\]
\medskip

\subsection*{``The System of the World''\\“世界体系”}

\noindent\textbf{English.} At the beginning of Book Three, which was devoted to ``The System of the World,'' Newton explained how it differed from the preceding two, which had been dealing with ``The Motion of Bodies'':
\par\noindent\textbf{中文。} 在第三卷开头，牛顿把该卷用于讨论“世界体系”，并说明它与前两卷的不同；前两卷讨论的是“物体运动”：
\medskip

\begin{quote}
\noindent\textbf{English.} In the preceding Books I have laid down \ldots{} principles not philosophical [pertaining to physics] but mathematical: such, namely, as we may build our reasonings upon in philosophical inquiries. These principles are laws and conditions of certain motions, and powers or forces, which chiefly have respect to philosophy; but, lest they should have appeared of themselves dry and barren, I have illustrated them here and there with some philosophical scholiums, giving an account of such things as are of a more general nature, and which philosophy seems chiefly to be founded on: such as the density and the resistance of bodies, spaces void of all bodies, and the motion of light and sounds. It remains that, from the same principles, I now demonstrate the structure of the System of the World.
\par\noindent\textbf{中文。} 在前几卷中，我提出了……并非哲学的[即属于物理学的]、而是数学的原理，也就是我们在哲学探究中可以据以推理的原理。这些原理是某些运动的定律和条件，以及能力或力；它们主要关系到哲学。不过，为免它们本身显得枯燥贫瘠，我在各处用一些哲学性释义加以说明，论述那些更具一般性、而哲学似乎主要建立于其上的事物，例如物体的密度和阻力、没有任何物体的空间，以及光和声音的运动。现在剩下的，就是从这些同一原理出发，证明世界体系的结构。
\end{quote}

\noindent\textbf{English.} I believe it fair to say that it was the freedom to consider problems either in a purely mathematical way or in a ``philosophical'' (or physical) way that enabled Newton to express the first law and to develop a complete inertial physics. After all, physics as a science may be developed in a mathematical way but it always must rest on experience--and experience never shows us pure inertial motion. Even in the limited examples of linear inertia discussed by Galileo, there was always some air friction and the motion ceased almost at once, as when a projectile strikes the ground. In the whole range of physics explored by Galileo there is no example of a physical object that has even a component of pure inertial motion for more than a very short time. It was perhaps for this reason that Galileo never framed a general law of inertia. He was too much a physicist.
\par\noindent\textbf{中文。} 我认为可以公允地说，正是这种自由——既可以用纯数学方式，也可以用“哲学的”（或物理的）方式来思考问题——使牛顿能够表述第一定律，并发展出完整的惯性物理学。毕竟，物理学作为一门科学可以以数学方式发展，但它始终必须建立在经验之上；而经验从未向我们显示纯粹的惯性运动。即使在伽利略讨论的线性惯性的有限例子中，也总有一些空气摩擦，而且运动几乎立刻终止，例如抛射体撞到地面时。在伽利略所探索的整个物理范围内，没有任何物理物体能在超过很短时间内具有哪怕一个纯粹惯性运动的分量。也许正因为如此，伽利略从未形成一条普遍的惯性定律。他太像一位物理学家了。
\medskip

\noindent\textbf{English.} But as a mathematician Newton could easily conceive of a body's moving along a straight line at constant speed forever. The concept ``forever,'' which implies an infinite universe, held no terror for him. Observe that his statement of the law of inertia, that it is the natural condition for bodies to move in straight lines at a constant speed, occurs in Book One of the \emph{Principia}, the portion said by him to be mathematical rather than physical. Now, if it is the natural condition of motion for bodies to move uniformly in straight lines, then this kind of inertial motion must characterize the planets. The planets, however, do not move in straight lines, but rather along ellipses. Using a kind of Galilean approach to this single problem, Newton could say that the planets must therefore be subject to two motions: one inertial (along a straight line at constant speed) and one always at right angles to that straight line, drawing each planet toward its orbit. (See, further, Supplements 11 and 12.)
\par\noindent\textbf{中文。} 但是作为数学家，牛顿很容易设想一个物体以恒定速度沿直线永远运动。“永远”这个概念暗含无限宇宙，却并不令他畏惧。请注意，他对惯性定律的表述——物体以恒定速度沿直线运动是其自然状态——出现在《原理》第一卷中，而牛顿说这一卷是数学的，而非物理的。现在，如果物体的自然运动状态是沿直线匀速运动，那么这种惯性运动必定也是行星的特征。然而，行星并不沿直线运动，而是沿椭圆运动。牛顿用一种类似伽利略的方法处理这个单一问题，可以说行星因此必定受制于两种运动：一种是惯性运动（以恒定速度沿直线运动），另一种总是与这条直线成直角，把每颗行星拉向其轨道。（另见补编11和12。）
\medskip

\noindent\textbf{English.} Though not moving in a straight line, each planet nevertheless represents the best example of inertial motion observable in the universe. Were it not for that component of inertial motion, the force that continually draws the planet away from the straight line would draw the planet in toward the sun until the two bodies collided. Newton once used this argument to prove the existence of God. If the planets had not received a push to give them an inertial (or tangential) component of motion, he said, the solar attractive force would not draw them into an orbit but instead would move each planet in a straight line toward the sun itself. Hence the universe could not be explained in terms of matter alone.
\par\noindent\textbf{中文。} 尽管行星并不沿直线运动，每颗行星仍然代表着宇宙中可观测到的惯性运动的最佳例子。如果没有那一惯性运动分量，那么持续把行星从直线拉开的力，就会把行星向太阳拉去，直到两个物体相撞。牛顿曾用这个论证来证明上帝的存在。他说，如果行星没有受到一次推动，从而获得惯性（或切向）运动分量，那么太阳的吸引力不会把它们带入轨道，而会使每颗行星沿直线朝太阳本身运动。因此，宇宙不能仅仅用物质来解释。
\medskip

\noindent\textbf{English.} For Galileo pure circular motion could still be inertial, as in the example of an object on or near the surface of the earth. But for Newton pure circular motion was not inertial; it was accelerated and required a force for its continuance. Thus it was Newton who finally shattered the bonds of ``circularity'' which still had held Galileo in thrall. And so we may understand that it was Newton who showed how to build a celestial mechanics based on the laws of motion, since the elliptical (or almost circular) orbital motion of planets is not purely inertial, but requires additionally the constant action of a force, which turns out to be the force of universal gravitation.
\par\noindent\textbf{中文。} 对伽利略来说，纯圆周运动仍然可以是惯性的，例如在地球表面或其附近的物体就是如此。但对牛顿来说，纯圆周运动不是惯性的；它是加速运动，并需要一个力来维持。因而，正是牛顿最终打碎了仍束缚着伽利略的“圆形性”枷锁。于是我们可以理解，正是牛顿展示了如何在运动定律基础上建立天体力学，因为行星的椭圆（或近乎圆形）轨道运动并非纯粹惯性运动，还需要一个力的持续作用；这个力最终证明就是万有引力。
\medskip

\noindent\textbf{English.} Thus Newton, again unlike Galileo, set out to ``demonstrate the structure of the System of the World,'' or--as we would say today--to show how the general laws of terrestrial motion may be applied to the planets and to their satellites.
\par\noindent\textbf{中文。} 因此，牛顿再次不同于伽利略，他着手“证明世界体系的结构”，或者用今天的话说，展示地上运动的一般定律如何能够应用于行星及其卫星。
\medskip

\begin{center}
\ldots
\end{center}

\noindent\textbf{English.} Isaac Newton's system of mechanics came to symbolize the rational order of the world, functioning under the ``rule of nature.'' Not only could Newtonian science account for present and past phenomena; the principles could be applied to the prediction of future events. In the \emph{Principia} Newton proved that comets are like the planets, moving in great orbits that must (according to Newtonian rules) be conic sections. Some comets move in ellipses and these must return periodically from far out in space to the visible regions of our solar system, whereas others will visit our solar system and never return. Edmond Halley applied these Newtonian results to an analysis of cometary records of the past and found--among others--a comet with a period of some seventy-five and a half years. He made a bold Newtonian prediction that this comet would reappear in 1758. When it did so, right on schedule, though long after Halley and Newton were dead, men and women everywhere experienced a new feeling of awe for the powers of human reason abetted by mathematics. This new respect for science was expressed by such adjectives as ``amazing,'' ``phenomenal,'' or ``extraordinary.'' This successful prediction of a future event symbolized the force of the new science: the perfection of the mathematical understanding of nature, realized in the ability to make reliable predictions of the future. Not surprisingly, men and women everywhere saw a promise that all of human knowledge and the regulation of human affairs would yield to a similar rational system of deduction and mathematical inference coupled with experiment and critical observation. The eighteenth century not only was the Enlightenment, but became ``preeminently the age of faith in science.'' Newton became the symbol of successful science, the ideal for all thought--in philosophy, psychology, government, and the science of society.
\par\noindent\textbf{中文。} 艾萨克·牛顿的力学体系逐渐成为世界理性秩序的象征，这一秩序在“自然法则”支配下运行。牛顿科学不仅能够解释现在和过去的现象；其原理还可以用于预测未来事件。在《原理》中，牛顿证明彗星与行星相似，它们沿巨大的轨道运动，而这些轨道按牛顿规则必定是圆锥曲线。有些彗星沿椭圆运动，必定会从遥远太空周期性返回太阳系的可见区域；另一些则会造访我们的太阳系后永不返回。埃德蒙·哈雷把这些牛顿结果应用于对过去彗星记录的分析，并在其中发现了一颗周期约为七十五年半的彗星。他作出一个大胆的牛顿式预言：这颗彗星将在1758年重现。后来它果然准时重现，尽管那时哈雷和牛顿都已去世很久，世界各地的人们仍对借助数学的人类理性力量产生了一种新的敬畏。这种对科学的新尊重，体现在“惊人”“非凡”或“卓越”等形容词中。对未来事件的成功预测，象征着新科学的力量：对自然的数学理解已臻完善，并表现为可靠预测未来的能力。不足为奇的是，世界各地的人们都看到了一种希望：人类全部知识以及人类事务的管理，都将服从一种类似的理性体系，这一体系把演绎和数学推理同实验及批判性观察结合起来。18世纪不仅是启蒙时代，而且成了“最突出地信仰科学的时代”。牛顿成为成功科学的象征，成为一切思想的理想典范——无论在哲学、心理学、政府，还是社会科学中。
\medskip

\noindent\textbf{English.} Newton's genius enables us to see the full significance of both Galilean mechanics and Kepler's laws of planetary motion as manifested in the development of the inertial principles required for the Copernican-Keplerian universe. A great French mathematician, Joseph Louis Lagrange (1736--1813), best defined Newton's achievement. There is only one law of the universe, he said, and Newton discovered it. Newton did not develop modern dynamics all by himself but depended heavily on certain of his predecessors; the debt in no way lessens the magnitude of his achievement. It only emphasizes the importance of such men as Galileo and Kepler, and Descartes, Hooke, and Huygens, who were great enough to make significant contributions to the Newtonian enterprise. Above all, we may see in Newton's work the degree to which science is a collective and a cumulative activity and we may find in it the measure of the influence of an individual genius on the future of a cooperative scientific effort. In Newton's achievement, we see how science advances by heroic exercises of the imagination, rather than by patient collecting and sorting of myriads of individual facts. Who, after studying Newton's magnificent contribution to thought, could deny that pure science exemplifies the creative accomplishment of the human spirit at its pinnacle?
\par\noindent\textbf{中文。} 牛顿的天才使我们能够看出伽利略力学和开普勒行星运动定律的全部意义，因为它们体现在哥白尼--开普勒宇宙所需要的惯性原理的发展之中。伟大的法国数学家约瑟夫·路易·拉格朗日（1736--1813）最恰当地概括了牛顿的成就。他说，宇宙只有一条定律，而牛顿发现了它。牛顿并不是独自发展出现代动力学的，他在很大程度上依赖某些前辈；这种债务丝毫不减损其成就的伟大。它只是强调了伽利略、开普勒、笛卡尔、胡克和惠更斯等人的重要性，他们足够伟大，能够为牛顿事业作出重要贡献。最重要的是，我们可以在牛顿的工作中看到，科学在多大程度上是一种集体的、累积的活动；也可以从中衡量一个个人天才对合作性科学努力之未来的影响。在牛顿的成就中，我们看到科学并不是靠耐心收集和分类无数个别事实而前进，而是靠想象力的英雄式运用而前进。研究过牛顿对思想的宏伟贡献之后，谁还能否认纯科学体现了人类精神创造性成就的顶峰？
\medskip

\section*{Text 3b: From \emph{The Principia: Mathematical Principles of Natural Philosophy}\\文本3b：选自《自然哲学的数学原理》}
\begin{flushright}
Isaac Newton
\end{flushright}

\subsection*{Definitions\\定义}

\subsubsection*{Definition 1\\定义一}
\noindent\textbf{English.} Quantity of matter is a measure of matter that arises from its density and volume jointly.
\par\noindent\textbf{中文。} 物质的量（即质量）是由密度和体积共同决定的物质量度。
\medskip

\noindent\textbf{English.} If the density of air is doubled in a space that is also doubled, there is four times as much air, and there is six times as much if the space is tripled. The case is the same for snow and powders condensed by compression or liquefaction, and also for all bodies that are condensed in various ways by any causes whatsoever. For the present, I am not taking into account any medium, if there should be any, freely pervading the interstices between the parts of bodies. Furthermore, I mean this quantity whenever I use the term ``body'' or ``mass'' in the following pages. It can always be known from a body's weight, for--by making very accurate experiments with pendulums--I have found it to be proportional to the weight, as will be shown below.
\par\noindent\textbf{中文。} 如果在一个也扩大为两倍的空间中，空气密度加倍，那么空气的量就是四倍；若空间扩大为三倍，则空气的量就是六倍。被压缩或液化而凝聚的雪和粉末也是如此；一切由于任何原因以各种方式被凝聚的物体也是如此。目前我并不考虑是否存在某种介质，能够自由贯穿物体各部分之间的空隙。此外，在以下各页中，我使用“物体”或“质量”一词时，指的就是这个量。这个量总是可以由物体的重量得知；因为通过非常精确的摆实验，我发现它与重量成正比，这将在后文说明。
\medskip

\subsubsection*{Definition 2\\定义二}
\noindent\textbf{English.} Quantity of motion is a measure of motion that arises from the velocity and the quantity of matter jointly.
\par\noindent\textbf{中文。} 动量是由速度和物质的量（质量）共同决定的运动量度。
\medskip

\noindent\textbf{English.} The motion of a whole is the sum of the motions of the individual parts, and thus if a body is twice as large as another and has equal velocity there is twice as much motion, and if it has twice the velocity there is four times as much motion.
\par\noindent\textbf{中文。} 整体的运动是各个部分运动的总和；因此，如果一个物体是另一个物体的两倍大且速度相同，它就有两倍的动量；如果它又有两倍的速度，它就有四倍的动量。
\medskip

\subsubsection*{Definition 3\\定义三}
\noindent\textbf{English.} Inherent force of matter is the power of resisting by which every body, so far as it is able, perseveres in its state either of resting or of moving uniformly straight forward.
\par\noindent\textbf{中文。} 物质的固有力，是每个物体借以抵抗的能力；凭借这种能力，物体尽其所能保持其静止状态，或沿直线作匀速运动的状态。
\medskip

\noindent\textbf{English.} This force is always proportional to the body and does not differ in any way from the inertia of the mass except in the manner in which it is conceived. Because of the inertia of matter, every body is only with difficulty put out of its state either of resting or of moving. Consequently, inherent force may also be called by the very significant name of force of inertia. Moreover, a body exerts this force only during a change of its state, caused by another force impressed upon it, and this exercise of force is, depending on the viewpoint, both resistance and impetus: resistance insofar as the body, in order to maintain its state, strives against the impressed force, and impetus insofar as the same body, yielding only with difficulty to the force of a resisting obstacle, endeavors to change the state of that obstacle. Resistance is commonly attributed to resting bodies and impetus to moving bodies; but motion and rest, in the popular sense of the terms, are distinguished from each other only by point of view, and bodies commonly regarded as being at rest are not always truly at rest.
\par\noindent\textbf{中文。} 这种力总是与物体成正比；除了理解方式不同之外，它与质量的惯性没有任何差别。由于物质具有惯性，每个物体都只能很困难地脱离其静止状态或运动状态。因此，固有力也可以用一个非常有意义的名称来称呼，即惯性力。此外，只有当另一个施加在物体上的力导致其状态改变时，物体才发挥这种力；这种力的发挥，依观察角度而定，既是抵抗，也是冲力：说它是抵抗，是因为物体为了保持其状态而努力反抗所施加的力；说它是冲力，是因为同一物体只是在困难地屈从于阻碍物的力时，努力改变那个阻碍物的状态。抵抗通常归于静止物体，冲力通常归于运动物体；但按通俗意义来说，运动和静止仅仅因观察角度不同而彼此区别，那些通常被认为静止的物体并不总是真正静止。
\medskip

\subsubsection*{Definition 4\\定义四}
\noindent\textbf{English.} Impressed force is the action exerted on a body to change its state either of resting or of moving uniformly straight forward.
\par\noindent\textbf{中文。} 施加力，是作用于物体、使其静止状态或沿直线匀速运动状态发生改变的作用。
\medskip

\noindent\textbf{English.} This force consists solely in the action and does not remain in a body after the action has ceased. For a body perseveres in any new state solely by the force of inertia. Moreover, there are various sources of impressed force, such as percussion, pressure, or centripetal force.
\par\noindent\textbf{中文。} 这种力仅存在于作用之中；作用停止后，它并不留存在物体内。因为物体完全凭借惯性力而保持在任何新的状态中。此外，施加力有多种来源，例如撞击、压力或向心力。
\medskip

\subsubsection*{Definition 5\\定义五}
\noindent\textbf{English.} Centripetal force is the force by which bodies are drawn from all sides, are impelled, or in any way tend, toward some point as to a center.
\par\noindent\textbf{中文。} 向心力，是物体从各个方向被吸引、被推动，或以任何方式趋向某一点如同趋向中心的力。
\medskip

\noindent\textbf{English.} One force of this kind is gravity, by which bodies tend toward the center of the earth; another is magnetic force, by which iron seeks a lodestone; and yet another is that force, whatever it may be, by which the planets are continually drawn back from rectilinear motions and compelled to revolve in curved lines.
\par\noindent\textbf{中文。} 这种力的一种是重力，物体借此趋向地球中心；另一种是磁力，铁借此趋向磁石；还有一种是使行星不断从直线运动中被拉回并被迫沿曲线运行的力，无论这种力究竟是什么。
\medskip

\noindent\textbf{English.} A stone whirled in a sling endeavors to leave the hand that is whirling it, and by its endeavor it stretches the sling, doing so the more strongly the more swiftly it revolves; and as soon as it is released, it flies away. The force opposed to that endeavor, that is, the force by which the sling continually draws the stone back toward the hand and keeps it in an orbit, I call centripetal, since it is directed toward the hand as toward the center of an orbit. And the same applies to all bodies that are made to move in orbits. They all endeavor to recede from the centers of their orbits, and unless some force opposed to that endeavor is present, restraining them and keeping them in orbits and hence called by me centripetal, they will go off in straight lines with uniform motion. If a projectile were deprived of the force of gravity, it would not be deflected toward the earth but would go off in a straight line into the heavens and do so with uniform motion, provided that the resistance of the air were removed. The projectile, by its gravity, is drawn back from a rectilinear course and continually deflected toward the earth, and this is so to a greater or lesser degree in proportion to its gravity and its velocity of motion. The less its gravity in proportion to its quantity of matter, or the greater the velocity with which it is projected, the less it will deviate from a rectilinear course and the farther it will go. If a lead ball were projected with a given velocity along a horizontal line from the top of some mountain by the force of gunpowder and went in a curved line for a distance of two miles before falling to the earth, then the same ball projected with twice the velocity would go about twice as far and with ten times the velocity about ten times as far, provided that the resistance of the air were removed.
\par\noindent\textbf{中文。} 用投石索旋转的石块努力离开旋转它的手；由于这种努力，它拉紧投石索，而且旋转得越快，拉得越强；一旦被释放，它就飞走。与这种努力相反的力，也就是投石索不断把石块拉回手的方向并使它保持在轨道中的力，我称之为向心力，因为它指向手，如同指向轨道的中心。所有被迫在轨道中运动的物体也都是如此。它们都努力离开自身轨道的中心；如果没有某种与这种努力相反的力存在，没有某种约束它们、使它们保持在轨道中并因而被我称为向心的力，它们就会沿直线作匀速运动而离去。如果一个抛射体失去重力，而且空气阻力被消除，它就不会向地球偏转，而会沿直线匀速飞向天空。抛射体由于自身重力而从直线路径被拉回，并不断向地球偏转；偏转程度的大小，与其重力和运动速度有关。相对于其物质的量（质量）而言重力越小，或者抛出时速度越大，它偏离直线路径就越少，飞得就越远。假如一个铅球从某座山顶被火药之力以给定速度沿水平线发射，在落到地面前沿曲线前进了两英里，那么在去除空气阻力的条件下，同一个球以两倍速度发射就会走大约两倍远，以十倍速度发射就会走大约十倍远。
\medskip

\noindent\textbf{English.} And by increasing the velocity, the distance to which it would be projected could be increased at will and the curvature of the line that it would describe could be decreased, in such a way that it would finally fall at a distance of 10 or 30 or 90 degrees or even go around the whole earth or, lastly, go off into the heavens and continue indefinitely in this motion. And in the same way that a projectile could, by the force of gravity, be deflected into an orbit and go around the whole earth, so too the moon, whether by the force of gravity--if it has gravity--or by any other force by which it may be urged toward the earth, can always be drawn back toward the earth from a rectilinear course and deflected into its orbit; and without such a force the moon cannot be kept in its orbit. If this force were too small, it would not deflect the moon sufficiently from a rectilinear course; if it were too great, it would deflect the moon excessively and draw it down from its orbit toward the earth. In fact, it must be of just the right magnitude, and mathematicians have the task of finding the force by which a body can be kept exactly in any given orbit with a given velocity and, alternatively, to find the curvilinear path into which a body leaving any given place with a given velocity is deflected by a given force.
\par\noindent\textbf{中文。} 通过增加速度，它被抛射到的距离可以任意增加，它所描绘曲线的曲率可以减小，以至于它最终会在相距10度、30度或90度处落下，甚至绕整个地球运行；最后，它也可能飞向天空，并无限期地继续这种运动。正如抛射体可以因重力而被偏转进入轨道并环绕整个地球一样，月球也可以——无论是由于重力（如果它有重力），还是由于任何其他使它趋向地球的力——总是从直线路径被拉回地球方向，并被偏转进入自己的轨道；没有这样的力，月球就不能保持在其轨道中。如果这个力太小，它就不能使月球充分偏离直线路径；如果太大，它就会使月球过度偏转，并把月球从轨道上拉向地球。事实上，这个力必须大小正合适；数学家的任务就是找出能够使物体以给定速度恰好保持在任一给定轨道中的力，反过来也要找出物体从任一给定地点以给定速度出发时，在给定力作用下被偏转成的曲线路径。
\medskip

\noindent\textbf{English.} The quantity of centripetal force is of three kinds: absolute, accelerative, and motive.
\par\noindent\textbf{中文。} 向心力的量有三种：绝对量、加速量和动力量。
\medskip

\begin{center}
\ldots
\end{center}

\noindent\textbf{English.} Further, it is in this same sense that I call attractions and impulses accelerative and motive. Moreover, I use interchangeably and indiscriminately words signifying attraction, impulse, or any sort of propensity toward a center, considering these forces not from a physical but only from a mathematical point of view. Therefore, let the reader beware of thinking that by words of this kind I am anywhere defining a species or mode of action or a physical cause or reason, or that I am attributing forces in a true and physical sense to centers (which are mathematical points) if I happen to say that centers attract or that centers have forces.
\par\noindent\textbf{中文。} 此外，也正是在同一意义上，我把吸引和推动称为加速性的和动力性的。并且，我可以互换而不加区别地使用表示吸引、推动或任何趋向中心之倾向的词，因为我并不是从物理观点，而只是从数学观点来考察这些力。因此，读者切勿以为我用这类词是在任何地方定义某种作用的种类或方式，或某种物理原因或理由；也不要以为当我偶尔说中心会吸引或中心具有力时，我是在真实的物理意义上把力归于中心（中心只是数学点）。
\medskip

\subsection*{Axioms, or the Laws of Motion\\公理，或运动定律}

\subsubsection*{Law 1\\定律一}
\noindent\textbf{English.} Every body perseveres in its state of being at rest or of moving uniformly straight forward, except insofar as it is compelled to change its state by forces impressed.
\par\noindent\textbf{中文。} 每一个物体都会保持其静止状态，或沿直线作匀速运动的状态，除非它被所施加的力迫使改变这种状态。
\medskip

\noindent\textbf{English.} Projectiles persevere in their motions, except insofar as they are retarded by the resistance of the air and are impelled downward by the force of gravity. A spinning hoop, which has parts that by their cohesion continually draw one another back from rectilinear motions, does not cease to rotate, except insofar as it is retarded by the air. And larger bodies--planets and comets--preserve for a longer time both their progressive and their circular motions, which take place in spaces having less resistance.
\par\noindent\textbf{中文。} 抛射体会保持其运动，除非它们受到空气阻力而迟滞，或被重力向下推动。一个旋转的环，其各部分由于凝聚力而不断把彼此从直线运动中拉回；除非受到空气迟滞，它不会停止旋转。更大的物体——行星和彗星——则在阻力较小的空间中，更长久地保持其前进运动和圆周运动。
\medskip

\subsubsection*{Law 2\\定律二}
\noindent\textbf{English.} A change in motion is proportional to the motive force impressed and takes place along the straight line in which that force is impressed.
\par\noindent\textbf{中文。} 动量的改变与所施加的动力成正比，并沿着施加该力的直线发生。
\medskip

\noindent\textbf{English.} If some force generates any motion, twice the force will generate twice the motion, and three times the force will generate three times the motion, whether the force is impressed all at once or successively by degrees. And if the body was previously moving, the new motion (since motion is always in the same direction as the generative force) is added to the original motion if that motion was in the same direction or is subtracted from the original motion if it was in the opposite direction or, if it was in an oblique direction, is combined obliquely and compounded with it according to the directions of both motions.
\par\noindent\textbf{中文。} 如果某个力产生某种动量，两倍的力就会产生两倍的动量，三倍的力就会产生三倍的动量；无论该力是一次施加，还是逐渐连续施加，都是如此。如果物体先前已经在运动，新的动量（由于动量总是与产生它的力同方向）若与原动量同方向，就加到原动量上；若与原动量相反，就从原动量中减去；若成斜向，则按照两个动量的方向作斜向结合并合成。
\medskip

\subsubsection*{Law 3\\定律三}
\noindent\textbf{English.} To any action there is always an opposite and equal reaction; in other words, the actions of two bodies upon each other are always equal and always opposite in direction.
\par\noindent\textbf{中文。} 对于任何作用，总有一个相反且相等的反作用；换言之，两个物体相互之间的作用总是大小相等，方向相反。
\medskip

\noindent\textbf{English.} Whatever presses or draws something else is pressed or drawn just as much by it. If anyone presses a stone with a finger, the finger is also pressed by the stone. If a horse draws a stone tied to a rope, the horse will (so to speak) also be drawn back equally toward the stone, for the rope, stretched out at both ends, will urge the horse toward the stone and the stone toward the horse by one and the same endeavor to go slack and will impede the forward motion of the one as much as it promotes the forward motion of the other. If some body impinging upon another body changes the motion of that body in any way by its own force, then, by the force of the other body (because of the equality of their mutual pressure), it also will in turn undergo the same change in its own motion in the opposite direction. By means of these actions, equal changes occur in the motions, not in the velocities--that is, of course, if the bodies are not impeded by anything else. For the changes in velocities that likewise occur in opposite directions are inversely proportional to the bodies because the motions are changed equally. This law is valid also for attractions, as will be proved in the next scholium.
\par\noindent\textbf{中文。} 凡是压迫或牵引别物者，也被别物以同样大小压迫或牵引。如果有人用手指压一块石头，手指也被石头压。如果一匹马牵引一块系在绳上的石头，那么这匹马（可以这样说）也同样被向石头方向拉回；因为绳子在两端被拉紧时，会由于同一个趋于松弛的努力，把马推向石头，也把石头推向马，并且它阻碍一方前进运动的程度，正如它促进另一方前进运动的程度一样。如果某个物体撞击另一个物体，并以自身的力以某种方式改变那个物体的动量，那么由于另一物体的力（因为二者相互压力相等），它自身也会反过来在相反方向上经历同样的动量改变。通过这些作用，发生相等的动量改变，而不是相等的速度改变——当然，前提是这些物体不受其他事物阻碍。因为既然动量改变相等，那么同样在相反方向发生的速度改变就与物体的质量成反比。这条定律也适用于吸引作用，这将在下一条释义中得到证明。
\medskip

\subsubsection*{Corollary 1\\推论一}
\noindent\textbf{English.} A body acted on by [two] forces acting jointly describes the diagonal of a parallelogram in the same time in which it would describe the sides if the forces were acting separately.
\par\noindent\textbf{中文。} 一个物体在[两个]力共同作用下所描绘的平行四边形对角线，所用时间与这些力分别作用时它描绘各边所用的时间相同。
\medskip

\begin{center}
\includegraphics[width=0.72\linewidth]{chunks/images/86e81927fb72b006792bc4605366d6baa2c5963aedd08e9e51d23ec99fb2611f.jpg}
\end{center}

\noindent\textbf{English.} Let a body in a given time, by force $M$ alone impressed in $A$, be carried with uniform motion from $A$ to $B$, and, by force $N$ alone impressed in the same place, be carried from $A$ to $C$; then complete the parallelogram $ABDC$, and by both forces the body will be carried in the same time along the diagonal from $A$ to $D$. For, since force $N$ acts along the line $AC$ parallel to $BD$, this force, by law 2, will make no change at all in the velocity toward the line $BD$ which is generated by the other force. Therefore, the body will reach the line $BD$ in the same time whether force $N$ is impressed or not, and so at the end of that time will be found somewhere on the line $BD$. By the same argument, at the end of the same time it will be found somewhere on the line $CD$, and accordingly it is necessarily found at the intersection $D$ of both lines. And, by law 1, it will go with [uniform] rectilinear motion from $A$ to $D$.
\par\noindent\textbf{中文。} 设一个物体在给定时间内，若只受在$A$处施加的力$M$作用，便以匀速运动从$A$被带到$B$；若只受同一地点施加的力$N$作用，便从$A$被带到$C$。现在补成平行四边形$ABDC$，则在两个力共同作用下，物体将在同一时间内沿对角线从$A$被带到$D$。因为力$N$沿$AC$线作用，而$AC$平行于$BD$；根据定律二，这个力不会使另一个力所产生的、朝向$BD$线的速度发生任何改变。因此，无论是否施加力$N$，物体都会在同一时间内到达$BD$线，所以在该时间结束时，它将位于$BD$线上的某处。用同样的论证，在同一时间结束时，它也将位于$CD$线上的某处；于是它必定位于两条线的交点$D$。并且，根据定律一，它将以[匀速]直线运动从$A$到$D$。
\medskip

\begin{center}
\ldots
\end{center}

\begin{center}

\includegraphics[width=0.9\linewidth]{chunks/images/96a7884762717ddc050668629e01a03990e49f38be99f72b422c072fe378ebb4.jpg}
\end{center}

\begin{center}
\includegraphics[width=0.9\linewidth]{chunks/images/453ff7a961a91fd9458260fc2cb186d44f83f69bb271740c7c73792307573ad1.jpg}
\end{center}

\begin{center}

\includegraphics[width=0.9\linewidth]{chunks/images/a57a6587536014598149960fe6cfc7bd98e69910545d5bb0e3e380f76308e9ba.jpg}
\end{center}
